\chapter{Introduction}
\label{ch:intro}

\section{Problem Statement}
\label{sec:problem}

Automated program repair (APR) has converged on a loop with two halves. A
\emph{proposer} produces a candidate patch, and an \emph{oracle} decides
whether that candidate is wrong by running it against a pool of test cases and
looking for one on which it misbehaves~\citep{legoues2019apr}. In the
generate-and-validate tradition the proposer was a mutation
operator~\citep{legoues2012genprog}; today it is usually a large language
model~\citep{xia2024chatrepair}. The oracle has not changed: it is still a
sequence of program executions, and it is still the half that burns the
machine.

How much it burns is easy to under-estimate, because the usual unit of account
hides it. The test-prioritisation literature reports \emph{test executions}
rather than time, and does so deliberately: the largest recent study of patch
validation states that it presents results in executions because ``the number
of test executions is more stable and more suitable for future work to
reproduce''~\citep{lou2024aprregression}. That choice is sound when executions
are interchangeable. On the corpus studied in this proposal they are not. The
oracle's own reference program takes a median of 0.076\,s of CPU on
one test case and 4.48\,s at the 99th percentile --- a $58.7\times$
spread --- and the variation is not between easy and hard \emph{tasks} but
inside a single task's pool, where the median coefficient of variation of
per-case cost is $0.571$. Counting executions on such a pool is counting apples
and buses.

The consequence shows up directly in the loop. A repair loop that walks the
pool in the order the benchmark happens to ship --- which is what any
implementation gets for free, and therefore what most implementations do ---
spends 2.18\,s of oracle CPU per round and keeps spending it: the
cost per round is flat from round~1 to round~20, because nothing in the loop
learns anything about the pool. Over the $6{,}800$ rounds measured in
Section~\ref{sec:pilot} that order costs 14843\,s of CPU where a
scheduled order costs 3532\,s.

Two kinds of memory have been proposed to make the loop cheaper, and neither
addresses this. \emph{Fault-recorded testing prioritisation} (\frtp)
ranks a test case by the number of patches it has already
falsified~\citep{qi2013frtp}; the same quantity, under the name ``patch killing
count'', drives later variants~\citep{venugopal2020mpp}. \cegmem, the study
that precedes this one, went further and stored every refuted patch under a
\emph{failure type} formed from the patch's edit location and the property it
violated, so that a fresh candidate matching a stored type could be rejected
before the oracle was paid at all. That worked: it cut oracle calls per
episode from $11.11$ to $1.57$, a factor of $7.07$
(Section~\ref{sec:cegmem-premise}).

But both memories are keyed by a \emph{class}, and a class-keyed memory answers
only one of the two questions a repair loop asks. It can answer \emph{which
patches are not worth testing}. It cannot answer \emph{which test to run
first}, and it certainly cannot answer \emph{which test is cheap} --- \frtp's
whole family records that a case killed, never what the kill cost, and never
the evidence of a case failing to kill, which is the information that
distinguishes a case that spared ten patches from one that was never run. In
\cegmem\ the residual was plainly visible: the memory removed most oracle
\emph{calls}, yet the loop still performed $39.5$ sandbox executions per
episode and still spent 14.1\,s per episode inside the guard,
because every surviving call still walked the pool in alphabetical order.

This proposal targets that residual. The idea is to key the memory by
\textbf{test case} instead of by failure class, to learn two online quantities
per case --- a kill rate and a cost --- and to schedule probes by their ratio.
Keying on the case has a second consequence that motivates the rest of the
design: a test case is an \emph{observation}, and only an observation can be
handed back to the proposer as a positive requirement (``produce exactly this
output on this input'') rather than as a prohibition (``do not do that
again''). The scheduler and the prompt contract are therefore one system, and
they ship together.

\section{Research Objectives and Questions}
\label{sec:rq}

This research aims to design, implement and evaluate a repair loop whose
memory of refuted attempts is indexed by test case, and to establish how much
of the oracle's cost that indexing removes. Four objectives structure the
work; each is paired with the question it answers and, following the discipline
adopted throughout this project, with the outcome that would \emph{refute} it.
A question with no refuting outcome is a name for a group of experiments, not a
question.

\begin{enumerate}[leftmargin=*,itemsep=6pt]

\item \textbf{Cost-aware probe scheduling.} Schedule the oracle's probes by an
online per-case kill rate divided by an online per-case cost, rather than by
pool order or by a kill count.
\emph{Research question (RQ1).} Does scheduling by \km\ reduce the CPU spent
refuting a patch, measured against the incumbent pool order, against an
uninformed random control, and against the published prioritisers \frtp\ and
additional-greedy coverage?
\emph{Refuted when} \km\ is not cheaper than pool order beyond the cluster
bootstrap's interval; or when it is cheaper than pool order but not cheaper
than \frtp\ and coverage, in which case the contribution is early stopping,
which is not ours.
\emph{Status.} Answered in the affirmative by the pilot of
Section~\ref{sec:pilot}; to be re-established on the full corpus in Phase~2.

\item \textbf{Attribution of the effect.} Decompose the index into the three
things it adds to the incumbent --- having a memory at all, turning a count
into a rate, and dividing that rate by cost.
\emph{Research question (RQ2).} Which of the three terms produces the effect?
\emph{Refuted when} no single step clears the bootstrap interval, so the gain
is ``having any memory'' and the score is ornament; or when the cost division
alone carries it, in which case the kill rate is ornament.
\emph{Status.} Partially answered by the pilot's exponent sweep
(Section~\ref{sec:pilot-sweep}), which separates the kill term from the cost
term; the full stepwise ladder is Phase~2.

\item \textbf{Renormalising the patch family.} Return the scheduled
counterexample to the proposer as a positive requirement keyed to a test case.
\emph{Research question (RQ3).} Does doing so move the proposal distribution
off the mode it has already drawn --- fewer duplicate patches, more repairs at
the same budget --- where a conversation transcript and a prohibition do not?
\emph{Refuted when} the duplicate rate does not fall against independent
repeated sampling; or when it falls and the repair rate does not rise, in which
case the distribution moved somewhere no better.
\emph{Status.} Not addressed by the pilot. The baseline is set by published
measurement: independent repeated prompting has already been reported to repair
$23$ bugs against conversational repair's $18$ at an equal budget of $24$ model
calls, with $65\,\%$ of the conversational system's patches
duplicates~\citep{chen2025conversational}. Any steering arm must beat repeated
sampling, not the pool order.

\item \textbf{The disclosure dose.} Showing a counterexample tells the model
what to satisfy.
\emph{Research question (RQ4).} At what level of disclosure does the proposer
begin satisfying \emph{the shown probes} rather than repairing the fault?
\emph{Refuted when} the share of patches refuted only outside the disclosed set
is flat across the number of facts shown and across disclosure level, in which
case there is no dose and the axis does not exist.
\emph{Status.} Phase~2.

\end{enumerate}

\paragraph{Contribution positioning.} None of the individual ingredients is
new. Cost-cognizant prioritisation dates to \citet{elbaum2001cost}; the
additional-greedy family to \citet{rothermel2001prioritizing}; prioritisation
inside a repair loop to \citet{qi2013frtp}; and Beta-posterior estimates with
$\varepsilon$-greedy exploration are standard. The claim made here is at the
level of the \emph{key and the unit}. First, the memory is indexed by test
case, which is what lets one record serve both the scheduler and the prompt;
every memory in the surveyed repair literature is indexed by a class. Second,
the ranked quantity is a rate divided by a measured cost, where the published
repair-loop prioritiser ranks on a count and discards both the non-kills and
the cost. Third --- and this is the finding the pilot contributes rather than
the design --- on a pool with $58.7\times$ internal cost spread the two units
of account disagree in sign, so the deployment-facing claim cannot be read off
an execution count at all.

\section{Scope}
\label{sec:scope}

The study is bounded to the dimensions in Table~\ref{tab:scope}.

\begin{table}[H]
\centering
\caption{Project scope.}
\label{tab:scope}
\small
\begin{tabularx}{\textwidth}{@{}lX@{}}
\toprule
\textbf{Dimension} & \textbf{Specification}\\
\midrule
Task domain & Single-fault repair of competitive-programming submissions; one
faulty program, one contest test pool, stdout-compared verdicts.\\
Benchmark & ConDefects~\citep{wu2024condefects}, Python subset. Chosen because
every test input ships with the output the judge accepted, so a refutation is a
concrete input/output pair rather than a pass/fail bit.\\
Proposer & One OpenAI-compatible chat model per grid; \code{gpt-4o-mini} for
the data the pilot reuses. The proposer is not a variable of the study.\\
Oracle & One program, one input, hard timeout, CPU measured with
\code{getrusage(RUSAGE\_CHILDREN)}; probing stops at the first refutation.\\
Cost unit & CPU seconds as primary, test executions reported alongside.
Wall-clock time is never used as a cost, because it is contaminated by sibling
processes.\\
Core method & Per-case Beta-posterior kill rate $K$ over an online per-case
mean cost $\hat{m}$, ranked by \km, with $\varepsilon$-greedy exploration.\\
Evidence in this document & A replay and re-execution pilot on the frozen
\cegmem\ data: $68$ tasks, $2{,}523$ cases, $4{,}670$ patches,
$6{,}800$ measured rounds per policy.\\
Out of scope here & The \caster\ grid itself (22 arms), the steered prompt
contract, duplicate-rate and token measurements, multi-fault and multi-language
repair, and any claim about repair \emph{quality} beyond the soundness
invariant of Section~\ref{sec:soundness}.\\
\bottomrule
\end{tabularx}
\end{table}

\paragraph{Scope of completed work at the proposal stage.} The \caster\
implementation exists: ten command-line verbs, the module boundaries of
Section~\ref{sec:overview}, and $263$ unit tests that run without a network, a
model or the benchmark. What has \emph{not} happened is the experiment. No arm
of the grid has been run to completion on the full corpus, no steered arm has
been run at all, and therefore nothing in this document derived from the
\caster\ repository is reported as a result. The evidence is the pilot, and the
pilot runs on the previous project's data by design --- see
Section~\ref{sec:why-pilot-here} for why that is the stronger choice and not
merely the available one.

\section{Report Structure}

Chapter~\ref{ch:related} surveys validation-cost work along four families ---
generate-and-validate repair and its validation bill, regression test
prioritisation, repair-specific prioritisation, and the token economics of
LLM repair loops --- consolidates them into one comparison table, and states
the gaps. Chapter~\ref{ch:background} gives the theory: the two costs of the
repair loop, total versus additional greedy prioritisation, what \frtp\ records
and what it throws away, the typed memory of \cegmem\ that this work starts
from, the Beta-posterior estimator behind $K$, and the cost-accounting rules.
Chapter~\ref{ch:framework} describes the proposed system and reports the
preliminary validation in full. Chapter~\ref{ch:evaluation} is the evaluation
plan: metrics, the arm roster, the ablation ladder, and the two tiers of
baseline. Chapter~\ref{ch:summary} summarises, lists limitations honestly, and
gives the phase timeline.
